\documentclass{report}
\usepackage{amsmath,amssymb}
\setlength{\parindent}{0mm}
\setlength{\parskip}{1em}
\begin{document}
\begin{center}
\rule{6in}{1pt} \
{\large Amanda Bienz \\
{\bf Reducing Communication Costs in Parallel Algebraic Multigrid}}

202 N Race Street \\ Apt 312 \\ Urbana \\ IL 61801
\\
{\tt bienz2@illinois.edu}\\
Rob Falgout\\
William Gropp\\
Luke Olson\\
Jacob Schroder\end{center}

Coarse-level fill-in as a result of a Galerkin product often yields
increased complexity and large communication costs in algebraic multigrid
(AMG) methods. Hence, the parallel efficiency of AMG may be improved by
reducing this communication. The communication cost associated with a
solver is dependent on the number of network links that are occupied at a
given time. This link occupancy is equal to the product of the number of
packets sent and the distance of each packet. Hence, the communication
cost is composed of the number of messages sent, the size of messages
sent, and the number of links traversed by each packet.

The number of messages in AMG, as well as their average size, can be
minimized by eliminating unnecessary communication, for example through
the use of non-Galerkin coarse grids. While this method can greatly
reduce the number of packets sent, it does not take into account the
distance these packets must travel. These methods are further improved by
targeting the costly communication through the use of topology-aware
methods. Removing small messages that traverse few links has little
effect on per-iteration cost, yet can have a large impact on convergence.
Furthermore, large messages that must travel a long distance can have a
larger effect on communication cost than convergence rate,
and their removal can be beneficial to the overall performance of AMG. In
this talk we give an overview of the parallel efficiency in non-Galerkin
methods and present directions that utilize topology to maximize
efficiency.


\end{document}
